\documentclass[../readings.tex]{subfiles}

\begin{document}

\subsection{Stability and Condition Number}
\label{sub:stability-theory}

Sensitivity of problems vs.\ robustness of algorithms.

\textBox{%
Conditioning asks whether the mathematical problem is sensitive to small changes in the data. Stability asks whether the algorithm introduces only small changes to the problem. A good computation needs both: a well-conditioned problem and a stable algorithm.
}

\subsubsection{Residual, Forward Error, and Backward Error}

\addDef{Residual and Forward Error}{%
For a computed solution $\hat{\bx}$ to the linear system in Def~\ref{def:linear-system}:
\begin{align*}
    \br = \bb-\bA\hat{\bx}
    \qquad \text{and} \qquad
    \be = \hat{\bx}-\bx^*,
\end{align*}
where $\bx^*$ is the exact solution.
\begin{itemize}
    \item \textbf{Residual:}\enspace How well the computed answer satisfies the equations.
    \item \textbf{Forward error:}\enspace How close the computed answer is to the exact answer.
\end{itemize}
}
\label{def:residual-forward-error}

\addExmpl{%
\textbf{(Small residual, large error)}\enspace Let
\begin{align*}
    \bA=\begin{pmatrix}1&0\\0&10^{-8}\end{pmatrix},
    \qquad
    \bb=\begin{pmatrix}1\\10^{-8}\end{pmatrix}.
\end{align*}
The exact solution is $\bx^*=(1,1)^T$. The approximation $\hat{\bx}=(1,0)^T$ has residual
\begin{align*}
    \br=\bb-\bA\hat{\bx}=(0,10^{-8})^T,
\end{align*}
which is tiny in absolute norm, but the solution error is $(0,-1)^T$. The small singular value makes one direction hard to determine accurately.
}

\addExcr{%
\begin{enumerate}
    \item Return to Remark~\ref{rem:residual-vs-error}. Explain why the residual in Def~\ref{def:residual-forward-error} is computable but the forward error usually is not.
    \item Verify every number in the small-residual example above.
    \item Change the small diagonal entry from $10^{-8}$ to $10^{-12}$. How do the residual and forward error change?
\end{enumerate}
}
\label{ex:residual-forward-error}

\addThm{Residual-Error Bound}{%
Let $\bA$ be nonsingular, $\bx^*$ solve $\bA\bx=\bb$, and $\hat{\bx}$ have residual $\br=\bb-\bA\hat{\bx}$. If $\bb\neq\bzero$, then
\begin{align*}
    \frac{\|\hat{\bx}-\bx^*\|}{\|\bx^*\|}
    \leq
    \kappa(\bA)\frac{\|\br\|}{\|\bb\|}.
\end{align*}
Thus a small relative residual guarantees a small relative forward error only when the problem is well-conditioned.
}
\label{thm:residual-error-bound}

\addExcr{%
\begin{enumerate}
    \item Starting from $\br=\bb-\bA\hat{\bx}$ and $\bb=\bA\bx^*$, show that $\br=\bA(\bx^*-\hat{\bx})$.
    \item Use the previous step to prove $\|\hat{\bx}-\bx^*\|\leq \|\bA^{-1}\|\|\br\|$.
    \item Use $\|\bb\|\leq \|\bA\|\|\bx^*\|$ to finish the proof of Thm~\ref{thm:residual-error-bound}.
    \item Apply the bound to the small-residual example above. Is the bound sharp?
\end{enumerate}
}
\label{ex:residual-error-bound}

\subsubsection{Why Factorize instead of Invert?}

\addRemark{%
\textbf{(Explicit inverse rule)}\enspace Never compute $\bA^{-1}$ explicitly.
\begin{itemize}
    \item \textbf{Cost:}\enspace Computing $\bA^{-1}$ takes $3\times$ longer than LU ($\frac{8}{3}n^3$ vs.\ $\frac{2}{3}n^3$).
    \item \textbf{Stability:}\enspace Round-off in $\bA^{-1}$ propagates into the entire product $\bA^{-1}\bb$. Stable factorizations confine errors.
    \item \textbf{Efficiency:}\enspace LU allows solving for new $\bb$'s in $O(n^2)$ without storing a dense inverse.
\end{itemize}
}

\addExcr{%
\begin{enumerate}
    \item Use the LU factorization costs from the LU chapter to compare total flops for solving 10 systems via $\bA^{-1}$ vs.\ LU when $n=500$.
    \item Construct $\bA = \text{diag}(1, 1, 10^{-12})$. Compare residuals of \texttt{inv(A) @ b} and \texttt{solve(A, b)} using Def~\ref{def:residual-forward-error}.
    \item Explain the result using the distinction between residual and forward error from Ex~\ref{ex:residual-forward-error}.
\end{enumerate}
}
\label{ex:inverse-vs-factorization}

\subsubsection{Numerical Stability}

A property of the \textbf{algorithm}.

\addDef{Backward Stability}{%
An algorithm is backward stable if its computed $\hat{\bx}$ is the \textbf{exact solution to a nearby problem}:
$(\bA + \delta \bA)\hat{\bx} = \bb + \delta \bb$, with $\|\delta \bA\|/\|\bA\| = O(\varepsilon_{\text{mach}})$.
}
\label{def:backward-stability}

\addThm{Normwise Backward Error for Linear Systems}{%
For a computed vector $\hat{\bx}$ with residual $\br=\bb-\bA\hat{\bx}$, the normwise relative backward error for the linear system $\bA\bx=\bb$ is
\begin{align*}
    \eta(\hat{\bx})
    =
    \frac{\|\br\|}
    {\|\bA\|\|\hat{\bx}\|+\|\bb\|}.
\end{align*}
This is the smallest relative perturbation size, measured normwise in both $\bA$ and $\bb$, that makes $\hat{\bx}$ an exact solution of a nearby system.
}
\label{thm:linear-system-backward-error}

\addExcr{%
\begin{enumerate}
    \item Suppose $(\bA+\delta\bA)\hat{\bx}=\bb+\delta\bb$. Show that $\br=\delta\bb-\delta\bA\hat{\bx}$.
    \item Use the triangle inequality to prove the lower bound
    \[
        \|\br\|\leq \|\delta\bA\|\|\hat{\bx}\|+\|\delta\bb\|.
    \]
    \item Divide by $\|\bA\|\|\hat{\bx}\|+\|\bb\|$ to obtain the denominator in Thm~\ref{thm:linear-system-backward-error}.
    \item Construct a rank-one perturbation in the direction of $\br$ to see why the bound can be attained.
    \item Compare this computable backward error with the residual-error bound in Thm~\ref{thm:residual-error-bound}.
\end{enumerate}
}
\label{ex:linear-system-backward-error}

\addRemark{%
\textbf{GEPP Stability:}\enspace
Gaussian Elimination with Partial Pivoting is backward stable for virtually all practical matrices. The growth factor $\rho$ stays small, preventing exponential error amplification.
}

\addRemark{%
\textbf{(Stable algorithms)}\enspace Householder QR and Cholesky for SPD systems are backward stable under their usual assumptions. LU with partial pivoting is reliable in practice, though its worst-case growth factor can be large.
}

\addExcr{%
\begin{enumerate}
    \item Use Def~\ref{def:backward-stability} to explain why backward stability is a natural standard when input data are already uncertain.
    \item Explain why a small residual does not always guarantee a small forward error. Use Thm~\ref{thm:residual-error-bound}.
    \item Compute $\eta(\hat{\bx})$ from Thm~\ref{thm:linear-system-backward-error} for the small-residual example above.
    \item Histogram growth factors for 100 random matrices via \texttt{scipy.linalg.lu}.
\end{enumerate}
}
\label{ex:backward-stability}

\subsubsection{Condition Number}

A property of the \textbf{problem}.

\addDef{Condition Number}{%
$\kappa(\bA) = \|\bA\|\|\bA^{-1}\|$.
\begin{itemize}
    \item \textbf{Interpretation:}\enspace Worst-case error amplification. A perturbation of size $\varepsilon$ in $\bb$ can cause error up to $\kappa(\bA) \varepsilon$ in $\bx$.
    \item \textbf{Rule of Thumb:}\enspace You expect at most $16 - \log_{10}(\kappa(\bA))$ correct digits in double precision.
\end{itemize}
}
\label{def:condition-number}

\addPrf{%
Let $\bA\bx = \bb$ be the original system and $\bA(\bx+\delta\bx) = \bb+\delta\bb$ the perturbed system. Subtracting gives $\bA \delta\bx = \delta\bb$, so $\delta\bx = \bA^{-1} \delta\bb$.
Taking norms:
\begin{align*}
    \|\delta\bx\| \leq \|\bA^{-1}\| \|\delta\bb\|.
\end{align*}
Also, from $\bA\bx = \bb$, we have $\|\bb\| \leq \|\bA\| \|\bx\|$, which implies $1/\|\bx\| \leq \|\bA\|/\|\bb\|$.
Combining these inequalities for the relative error:
\begin{align*}
    \frac{\|\delta\bx\|}{\|\bx\|} &\leq \|\bA^{-1}\| \|\delta\bb\| \frac{\|\bA\|}{\|\bb\|} \\
    &= (\|\bA\| \|\bA^{-1}\|) \frac{\|\delta\bb\|}{\|\bb\|} \\
    &= \kappa(\bA) \frac{\|\delta\bb\|}{\|\bb\|}.
\end{align*}
Thus, the condition number $\kappa(\bA)$ acts as the amplification factor for relative perturbations in the right-hand side.
}

\addExcr{%
\begin{enumerate}
    \item Reproduce the perturbation argument above without looking.
    \item Compute $\kappa_2(H_n)$ for Hilbert matrices $n=3, \dots, 10$. How fast does it grow?
    \item Use Def~\ref{def:backward-stability}, Def~\ref{def:condition-number}, and Thm~\ref{thm:residual-error-bound} to explain why a stable algorithm on an ill-conditioned problem may still lose many digits.
    \item \textbf{The Neumann Series:}\enspace Use it to prove the perturbation bound: $\frac{\|\delta \bx\|}{\|\bx\|} \leq \frac{\kappa(\bA)}{1 - \dots}$.
\end{enumerate}
}
\label{ex:conditioning-and-stability}

\subsubsection{Solver Selection Hierarchy}

\addThm{Solver Selection by Structure}{%
Choose a solver from the most specific trustworthy structure in the problem:
\begin{enumerate}
    \item If $\bA$ is SPD and dense, use Cholesky; if it is SPD and large sparse, use CG with a preconditioner.
    \item If $\bA$ is general, square, dense, and nonsingular, use LU with partial pivoting.
    \item If the problem is full-rank least squares, use Householder QR.
    \item If the problem is rank-deficient, nearly rank-deficient, or requires numerical rank information, use the SVD.
    \item If $\bA$ is large, sparse, nonsymmetric, and only matrix-vector products are practical, use GMRES or another nonsymmetric Krylov method.
\end{enumerate}
The more structure a solver exploits, the cheaper it can be; the less trustworthy the structure, the more robust the solver must be.
}
\label{thm:solver-selection}

\addExcr{%
\begin{enumerate}
    \item For each case in Thm~\ref{thm:solver-selection}, identify the mathematical structure being used: symmetry, positive definiteness, sparsity, full column rank, or numerical rank.
    \item Explain why Cholesky is inappropriate if the SPD assumption fails.
    \item Explain why QR is preferred to normal equations for full-rank least squares, using Ex~\ref{ex:normal-equations-conditioning}.
    \item Explain why the SVD is the natural fallback when rank is uncertain, using Thm~\ref{thm:svd-coordinates}.
    \item Explain why CG and GMRES fit the Krylov approximation principle from Thm~\ref{thm:krylov-approximation-principle}.
\end{enumerate}
}
\label{ex:solver-selection-by-structure}

\addDef{Iterative Refinement}{%
The accuracy of a computed solution $\hat{\bx}$ can be improved by iterative refinement:
\begin{enumerate}
    \item Compute the residual $\br = \bb - \bA\hat{\bx}$.
    \item Solve $\bA \boldsymbol{\delta} = \br$.
    \item Update the estimate: $\hat{\bx} \leftarrow \hat{\bx} + \boldsymbol{\delta}$.
\end{enumerate}
Each step recovers digits lost due to ill-conditioning, up to the limits of machine precision.
}
\label{def:iterative-refinement}

\addRemark{%
\textbf{(Limitations)}\enspace Iterative refinement helps when the residual can be computed accurately and the correction equation can be solved reliably. It cannot overcome a problem whose condition number is so large that the desired digits are not present in the data.
}

\addExcr{%
\begin{enumerate}
    \item Solve $H_{10}\bx = H_{10}\mathbf{1}$. Apply one step of iterative refinement from Def~\ref{def:iterative-refinement}. How many digits are recovered?
    \item Compare forward errors of \texttt{solve} (LU) vs \texttt{lstsq} (QR) on a Hilbert system. Interpret the result using Ex~\ref{ex:conditioning-and-stability}.
    \item Choose a solver for each case in Thm~\ref{thm:solver-selection}: SPD, dense general square, full-rank least squares, rank-deficient least squares, large sparse SPD, and large sparse nonsymmetric. Justify each choice by citing one earlier result or exercise.
\end{enumerate}
}
\label{ex:solver-selection-refinement}

\end{document}
